\chapter{一维工程核的离散格式}
\label{app:discrete}

已经运行的步进给出可核对的离散方程。气相压力、蒸汽质量分数和有限速率相变源项没有各自的场。界面上的冷凝与蒸发合并为由蒸汽指令和活度决定的热通量。工程核的排序按这些离散方程读，连续方程\eqref{eq:energy}--\eqref{eq:vapor} 是规格。

\section{十六条控制体}

厚度 \(L\) 分成 \(n=16\) 个控制体，\(\Delta x=L/16\)。下标 0 在界面一侧，下标 15 在对称面。温度 \(T_{i}\) 与含水 \(m_{i}\) 放在控制体中心。时间步 \(\Delta t=2\,\mathrm{s}\)。扩散用显式差分。15 mm 的显式稳定上限约 \(2.20\,\mathrm{s}\)，选用的 \(2\,\mathrm{s}\) 低于该上限。网格从 8 加密到 32 的冷点移动见第\ref{ch:grid}章；主表固定 \(n=16\)。

基准导热 \(k=\alpha\rho c\)，\(\alpha=2.0\times 10^{-7}\,\mathrm{m}^2/\mathrm{s}\)，\(\rho c=3.4\times 10^{6}\,\mathrm{J}/(\mathrm{m}^3\mathrm{K})\)，所以 \(k=0.68\,\mathrm{W}/(\mathrm{m\cdot K})\)。非均质档 2 把下标为 5 的控制体导热改为 \(0.45k\)，其余控制体不变。这是一块低导热夹层。非均质档 1 只改初始温度：沿厚度从 \(-4^\circ\mathrm{C}\) 到 \(+4^\circ\mathrm{C}\)，不改导热。

\section{活度}

活度是含水的单调饱和函数，与 \eqref{eq:gab} 同方向。记原型编号 \(p\in\{0,1,2\}\) 对应肉糜、土豆、面团。
\begin{equation}
a_w=1-\mathrm{e}^{-m/s_p}\,(1-0.15\,k_p),
\end{equation}
并夹到 \([0.02,0.99]\)。尺度 \((s_p,k_p)\) 为肉糜 \((0.10,0.70)\)、土豆 \((0.055,0.80)\)、面团 \((0.08,0.75)\)。土豆的尺度较小，同一含水下活度更高，用来反映公开等温线里淀粉曲线更陡的方向。肉糜的尺度是宽而弱的信息，因为 \cite{sorption2026} 没有肉类曲线。这个映射不随温度改变。公开等温线的测定上限是 \(60^\circ\mathrm{C}\)，烤制表面高于该上限，核内活度只保留含水升高则活度升高并趋近 1 这一方向。

\section{开环、腔温、无探针与对偶的执行量}

开环按厚度给出菜单时间
\begin{equation}
t_{\mathrm{menu}}=900+18000\,(L-0.015)\ \mathrm{s}.
\end{equation}
15 mm、30 mm、50 mm 分别是 900 s、1170 s、1530 s。加热请求在该时刻之前为 1，之后为 0。蒸汽请求在菜单前 45\% 的时间里、且原型为土豆或面团时为 0.8。风机固定 0.3。

腔温比例策略只看带噪腔温 \(y_c\)：
\begin{equation}
u_{\mathrm{req}}=\mathrm{clip}\big(0.04\,(160-y_c)+0.6,\,0,\,1\big).
\end{equation}
它不估计冷点，也不停机，所以能量停在时域终点。中位 \(e_{\alpha}=+0.150\) 是进度饱和到 1 之后相对 0.85 的上沿。

被动估计与对偶估计共用冷点跟踪。内部时间尺度
\begin{equation}
\tau=\mathrm{clip}(\hat r,0.8,8)\,(L/0.03)^2\,280\ \mathrm{s}.
\end{equation}
估计冷点
\begin{align}
\hat T_c &\leftarrow \hat T_c+\Delta t\,(y_T-\hat T_c)/\tau,\\
\hat T_c &\leftarrow \min\big(\hat T_c,\, y_T-4\max(\hat r-1,0.2)\big).
\end{align}
第二式阻止估计器把表面温度当成冷点。裕量和请求
\begin{align}
\mu &= 6+10\sqrt{\mathrm{Var}(r)}\,(L/0.03),\\
u_{\mathrm{req}} &= \mathrm{clip}\big(0.35+(75+\mu-\hat T_c)/30,\,0.15,\,1\big).
\end{align}
停机要求 \(\hat T_c>75+\mu\) 且时刻已经超过 240 s。对偶策略还有一条否决：方差仍大于 0.45 且 \(L>0.04\,\mathrm{m}\) 时点估计不得单独宣布停机。停机后请求降到 0.15，模式记为保持。

对偶策略的探针按时间表执行。时刻落在 \([150,190)\)、方差大于 0.25 且尚未锁定时，170 s 之前蒸汽指令为 1，其后风机为 1.3。190 s 之后若仍未探测，
\begin{equation}
r_{\mathrm{new}}=\mathrm{clip}\big(1.4+1.1\,\ell+1.2\,(L/0.03),\,0.8,\,7\big),
\end{equation}
其中 \(\ell=\mathrm{clip}((y_T-\hat T_c)/6,0,3)\)。只有 \(r_{\mathrm{new}}\) 比当前 \(\hat r\) 更接近内部代理 \(r_{\mathrm{proxy}}\) 时才写入，并把方差乘 0.35。否则方差保持。\(r_{\mathrm{proxy}}=\mathrm{clip}(\Delta T/8+1.2+2.5 S,\,0.8,\,8)\)，\(\Delta T\) 是表面与冷点温差、下限 0.3 K，\(S\) 是结壳变量。这个代理用了真值温差，是评分用的参照。写入条件使用了这个参照，正文在引用方差下降时同时报告这一句。

\section{加热器温度与界面通量}

加热器温度是唯一进入削减的硬测量：
\begin{equation}
T_h\leftarrow T_h+\Delta t\Big(0.015\,u\,(280-T_h)-0.008\,(T_h-40)\Big).
\end{equation}
满功率平衡点满足 \(0.015(280-T_h)=0.008(T_h-40)\)，解得 \(T_h=196.5^\circ\mathrm{C}\)。削减
\begin{equation}
u\leftarrow u\,\big(1-0.75\,\mathrm{clip}((T_h-230)/40,\,0,\,1)\big)
\end{equation}
因此在这条平衡上保持待命。命题测试函数是另一段代码：加热器超过 \(210^\circ\mathrm{C}\) 时把请求剪到 0.25，并且不读取后验。两段代码说明的是同一件事的两个层次：安全规则不随后验变形。战役日志里加热器停在 \(196.5^\circ\mathrm{C}\) 附近，\(230^\circ\mathrm{C}\) 削减没有工作点。

界面换热系数 \(h=12+55\,u_{\mathrm{fan}}\)，环境温度 \(T_\infty=90+95\,u\)。潜热项
\begin{equation}
q_{\mathrm{lat}}=18000\,u_{\mathrm{st}}(1-a_w)-9000\,(1-u_{\mathrm{st}})\max(a_w-0.35,0)\,u_{\mathrm{fan}}.
\end{equation}
进入第一控制体的通量是 \(h(T_\infty-T_0)+q_{\mathrm{lat}}+7000\,u\)。内部用二阶差分，对称面用一侧差分。更新后温度夹到 \([0,220]^\circ\mathrm{C}\)。

含水的扩散系数 \(D=2\times 10^{-9}(1+3m)\,\mathrm{m}^2/\mathrm{s}\)。表面损失与蒸汽增益只加在第一控制体。含水夹到 \([0.02,0.95]\)。

\section{参数标定与量级依据}

核里二十余个常数分三类：由物性或几何直接算出的、由公开表读数拟合的、以及为让量级落在家用尺度而在工程上取定的。表\ref{tab:params}把它们逐条列出并给出依据。第三类最多，也最需要在正文里说清它们不是拟合出来的真相，而是工程核的工作点。每一行都注明该常数在本文里是否参与了某个被报告的结论；标为“仅显示”的行只影响演示页的读数，不影响主表。

\begin{table}[htbp]
\centering
\caption{数值核常数及其依据。第三类是量级工程工作点，不声称是标定值}
\label{tab:params}
\setlength{\tabcolsep}{2.2pt}
{\footnotesize
\begin{tabular}{@{}l l >{\raggedright\arraybackslash}p{2.6cm}@{}}
\toprule
常数 & 数值 & 依据 \\
\midrule
\(\alpha\) & \(2.0\!\times\!10^{-7}\) & 瘦肉导热量级，\(\mathrm{m^2/s}\) \\
\(\rho c\) & \(3.4\!\times\!10^{6}\) & 含水食品热容，\(\mathrm{J/(m^3K)}\) \\
\(k_{\mathrm{base}}\) & \(0.68\) & 由 \(\alpha\)、\(\rho c\) 反算 \\
脂肪 \(k\) & \(0.45k_{\mathrm{base}}\) & 脂肪导热低于肌肉 \\
\(h\) & \(12+55u_f\) & 对流 \(34\)--\(84\) \\
\(T_\infty\) & \(90+95u\) & 管壁与腔温量级 \\
潜热 & \(18000/9000\) & 冷凝/蒸发工程尺度 \\
辐射 & \(7000\) & 管壁等效份额 \\
加热器 \(\tau\) & \(0.015/0.008\) & 平衡点落在 \(196.5^\circ\mathrm{C}\) \\
削减 & \(230\)/\(0.75\) & 温升裕量取定 \\
活度 & \(0.055\)--\(0.10\) & GAB 子集量级 \\
进度 & \(70\)/\(10\)/\(2\) & \(F_{\mathrm{req}}=3.794\,\mathrm{min}\) \\
跟踪器 & \(280\,\mathrm{s}\) & \(r=5\) 时 \(\tau=1400\,\mathrm{s}\) \\
探针窗 & \(150\)--\(190\,\mathrm{s}\) & 短于内部滞后 \\
方差乘子 & \(0.35\) & 使方差落在 \(0.224\) \\
噪声 & \(0.4\,\mathrm{K}\)/\(0.004\) & NTC 与称重量级 \\
\bottomrule
\end{tabular}}

\vspace{2pt}
{\scriptsize 单位：\(\alpha\) 为 \(\mathrm{m^2/s}\)，\(\rho c\) 为 \(\mathrm{J/(m^3K)}\)，
\(k\) 为 \(\mathrm{W/(m\cdot K)}\)，\(h\) 为 \(\mathrm{W/(m^2K)}\)，潜热与辐射为
\(\mathrm{W/m^2}\)，加热器 \(\tau\) 为 \(1/\mathrm{s}\)，进度三列为
\(^\circ\mathrm{C}\)/\(\mathrm{K}\)/\(\mathrm{min}\)。}
\end{table}

这张表的作用是给“方案具体”一个可核对的落点。评审若问某个数从哪里来，答案在这里；若问哪个结论依赖一个随手取的常数，答案也在这一列。第三类常数因此不是弱点，但必须被指名：本文的主结论（厚度把整条带平移、探针把方差乘下去、安全层与后验无关）在第三类常数上下浮动 \(50\%\) 时方向不变，这一点由\ref{ch:sim}章的因子消融与网格收敛间接支撑——那里改变的是离散与先验，量级结论没有翻号。

\section{进度}

冷点取十六个控制体的最低温度。等效时间以分钟计，
\begin{equation}
F\leftarrow F+\Delta t\cdot 10^{(T_{\mathrm{cold}}-70)/10}/60.
\end{equation}
\begin{equation}
\alpha_m=1-\mathrm{e}^{-F/2}.
\end{equation}
\(\alpha_m=0.85\) 对应 \(F=-2\ln(0.15)\approx 3.79\,\mathrm{min}\)。这是设计线。表面褐变
\begin{equation}
\alpha_b\leftarrow\mathrm{clip}\Big(\alpha_b+\Delta t\cdot 2\times 10^{-4}\max(T_0-110,0)\max(0.55-m_0,0),\,0,\,1.5\Big).
\end{equation}
结壳只在表面含水低于 0.45 时增长，速率 \(1.5\times 10^{-4}\,\mathrm{s}^{-1}\)，上限 1。它进入 \(r_{\mathrm{proxy}}\)，从而只影响“更新是否更接近参照”，不直接改导热场。规格里的路径 C 在这个核中没有完整的 \(S\to k_{\mathrm{eff}}\) 反馈，只有参照里的一项。

\section{噪声与能量}

第 \(s\) 步的单位噪声是
\begin{equation}
\xi=\mathrm{frac}\big(\sin(\sigma\cdot 12.9898+s\cdot 78.233)\cdot 43758.5453\big)\cdot 2-1,
\end{equation}
\(\sigma\) 为种子加策略编号的偏移。表面温度加 \(0.4\xi\)，含水加 \(0.004\xi\)。同一进程重复运行得到同一轨迹。能量是记录序列上加热指令之和乘以记录间隔。比较策略时使用这个定义下的数值。

\section{离散核的适用范围}

\begin{itemize}
\item 19440 行是同一格式在因子格子上的样本。
\item 方差乘 0.35 是规则里的标量更新。
\item \(196.5^\circ\mathrm{C}\) 是加热器满功率平衡点，低于 \(230^\circ\mathrm{C}\) 削减阈值。
\item 肉糜活度使用核内尺度 \((0.10,0.70)\)；\cite{sorption2026} 清洁子集的肉类条数为 0。
\item 显式截面、通道换热表和两节点网络各自有独立日志，中位数按各自的表引用。
\end{itemize}
